\section{Exact Kuga--Satake correspondences}
\label{sec:exactks}

We now use the realization and deformation results to construct the
prescribed Kuga--Satake map.  The distinction between that map and an
unspecified nonzero embedding matters: the full Clifford action will
subsequently control tensor classes and the actions on different K3
factors.

For a projective K3 surface $S$, put
\[
 T(S)=\NS(S)_{\Q}^{\perp}\subset H^2(S,\Q),
 \qquad W_S=C^+(T(S),q),
\]
where $q$ is the cup form and our Clifford convention is
$v^2=q(v,v)$.  Give $W_S$ its standard left Kuga--Satake Hodge
structure.  Thus, for an oriented orthonormal basis $u_1,u_2$ of
the positive period plane, its complex structure is left
multiplication by $u_1u_2$.  This polarizable weight-one structure
is $H^1(A_S,\Q)$ for an abelian variety, up to the choice of lattice
and isogeny \cite{KugaSatake,DeligneK3,VanGeemenKS}.  If
$w\in T(S)$ has nonzero square, the standard tensor is
\begin{equation}
 \label{ks:standard-map}
 j_w:T(S)(1)\longrightarrow\End(W_S),
 \qquad j_w(v)(x)=vxw.
\end{equation}
A rational polarization identifies $W_S^\vee$ with $W_S(1)$ and
therefore converts \eqref{ks:standard-map} into a Hodge map
$\iota_w:T(S)\to W_S\otimes W_S\subset H^2(A_S^2,\Q)$.

\begin{theorem}
 \label{thm:exact-ks}
Let $S$ be a smooth projective complex K3 surface.  For every
nonisotropic $w\in T(S)$, every full even-Clifford realization
$H^1(A_S,\Q)=C^+(T(S),q)$, and every rational polarization used
to define $\iota_w$, there is a cycle
\[
 \Gamma\in\CH^2(S\times A_S^2)_{\Q}
 \quad\text{such that}\quad
 \Gamma_*|_{T(S)}=\iota_w.
\]
\end{theorem}

The proof first constructs this exact map at the mirror quartic
$B$.  Auxiliary terms in its Mukai vector make the cap-action
kernel rigid enough to determine the unknown transcendental
component.  Equality in the Ext bound then permits deformation
through a full polarized K3 family.  A rational comparison of
polarization lines and specialization reach every projective K3
surface.

\subsection{A spin-invariant tensor at the mirror quartic}
\label{ks:construction}

Retain the spaces $P_B$ and $U$ from the cohomological dictionary.
Choose a rational ample class $h\in\NS(B)_{\Q}$ and set
\[
 N=h^\perp\cap\NS(B)_{\Q},\qquad V=N\perp U.
\]
Then $\dim N=18$, while $V$ has dimension $21$ and signature
$(2,19)$.  The square class of $q(h,h)$ can be prescribed as
$2d$ for any positive integer $d$; we explain this choice when
comparing polarization degrees in Section~\ref{ks:deformation}.

For this $V$, write
\[
 W=C^+(V).
\]
The left spin action on $W$ admits a nondegenerate invariant
alternating form that polarizes the Kuga--Satake structure on
the chosen period component.  One may use
\begin{equation}
 \label{ks:trace-form}
 E(x,y)=\operatorname{Tr}(\alpha x^*y),
\end{equation}
where $*$ reverses Clifford products and $\alpha^*=-\alpha$.
For example, take $\alpha$, up to sign, to be the product of
two orthogonal rational vectors spanning a positive plane.
At that plane, let $I$ be their normalized product.  Evaluating
$\operatorname{Tr}(I x^* I x)$ on the even Clifford monomial
basis gives a definite form: commuting the two positive
generators past a monomial compensates the parity of its
negative generators.  Choose the sign giving positivity.
Spin invariance then gives positivity on the entire component.

Let $\theta\in\bigwedge^2W$ be the nondegenerate alternating
tensor corresponding to this form, and choose a rational volume
element $z$ in the full Clifford algebra of $V$.  Define
\begin{equation}
 \label{ks:J}
 J:V\longrightarrow\bigwedge^2W,
 \qquad J(v)=(L_{vz}\otimes1)\theta.
\end{equation}
Here $L_a$ denotes left multiplication by $a$.  Since $\dim V=21$,
the volume element is central and $vz$ is reversal-even.
Consequently $L_{vz}$ is self-adjoint for \eqref{ks:trace-form},
so \eqref{ks:J} is alternating.  It is also Spin-equivariant.

Restriction to the split group $\Spin(U)=\mathrm{SL}_2$ gives
\[
 W=M_0\otimes H,
 \qquad H=\Q e_+\oplus\Q e_-.
\]
Thus this is precisely a torus space in the preceding dictionary,
with $g=\dim M_0=2^{19}$.  To describe the operators on this large
multiplicity space, work over $\C$.  The even Clifford algebra in
dimension $21$ is the endomorphism algebra of its spin module
$\Sigma$, of dimension $2^{10}$.  Its regular representation is
therefore a sum of $2^{10}$ copies of $\Sigma$.  On restriction to
$N\perp U$, the module $\Sigma$ is $M_1\otimes H$, where $M_1$
is the $2^9$-dimensional irreducible module for the full Clifford
algebra of $N_{\C}$.  Consequently
\begin{equation}
 \label{ks:multiplicity}
 W_{\C}=\Sigma\otimes R_0=(M_1\otimes H)\otimes R_0,
 \qquad M_{0,\C}=M_1\otimes R_0,
\end{equation}
where $R_0$ is the regular-representation multiplicity space.
The module $\Sigma$ has an invariant alternating form.  The tensor $\theta$ is the
product of a symmetric nondegenerate tensor $Q$ on $M_0$ with
the alternating tensor on $H$.  Moreover, $Q$ factors as the
form on $M_1$ times a symmetric form on $R_0$.

For $y\in N$, the operator $L_{yz}$ is, up to a common nonzero
scale, the Clifford generator $\Gamma_y$ on $M_1$ tensored with
the identity on $R_0\otimes H$.  These generators are
self-adjoint for $Q$.  Let $\Delta$ be their chirality operator;
it anticommutes with every $\Gamma_y$ and is skew-adjoint,
since reversal in dimension $18$ has negative sign.  For
$u\in U$, $L_{uz}$ is, up to a common scale, $\Delta$ times
the vector action on $H$.  It follows that $J(y)$ has symmetric
matrix proportional to $\Gamma_yQ$ and alternating $H$ factor,
whereas $J(u)$ has matrix proportional to $\Delta Q$ and
belongs to the $\operatorname{Sym}^2H$ part.

The Clifford generators on $M_1$ generate its full endomorphism
algebra.  An operator commuting with all of them can still act
arbitrarily on $R_0$.  We remove this multiplicity freedom by choosing
two further rational Spin-invariant nondegenerate
alternating tensors $\theta_1,\theta_2$ whose relative
operators to $\theta$ have scalar common commutant on $R_0$.
Such choices exist: over $\C$ two general symmetric operators
relative to the multiplicity form have this property, which is
a nonempty Zariski-open condition in the rational space of
invariant tensors.

We next place these tensors in separate exterior degrees.  This will
let an annihilation equation distinguish the wedge, derivation, and
contraction parts of a torus operator.  Choose a rational basis
$e_1,e_2,e_3$ of $N^\perp\cap P_B$,
each vector lying in one of $\Q h$, $H^0(B,\Q)$, and
$H^4(B,\Q)$.  For dual bases $(v_j)$ and $(v_j^\vee)$ of
$V$, put
\begin{equation}
 \label{ks:test-tensor}
 \begin{split}
 R_1&=1,\qquad
 R_2=\theta^4+\theta_1^7+\theta_2^{10},\qquad
 R_3=\theta^{13},\\
 C_*&=\sum_jv_j\otimes J(v_j^\vee)
            +\sum_{s=1}^3e_s\otimes R_s.
 \end{split}
\end{equation}
The first summand contains the desired Clifford map $J$.
The three $e_s$ components impose the additional equations just
described; after those equations remove the unwanted operators,
the remaining kernel will consist of the $19$ simultaneous spin
lowering directions.  We prove this next.
The full tensor is invariant under the simultaneous $\mathrm{SL}_2$
action.  Its inverse image under the dictionary is therefore
a rational middle-degree class $a$ killed by the symplectic
form.

\subsection{The formal kernel has dimension nineteen}
\label{ks:formal-kernel-section}

For a tensor in $U\otimes\bigwedge W$, write its three weight
components as
\begin{equation}
 \label{ks:weight-components}
 t_+\otimes U_-+t_0\otimes U_0+t_-\otimes U_+.
\end{equation}
Here $U_-,U_0,U_+\in\bigwedge W$.  A superscript $*$ denotes
the components of $C_*$; we will later use the same notation for
the component in $U_B\otimes\bigwedge W$.
In particular, $U_-^*$ is nonzero and has exterior degree two.
The annihilator of $a$ by degree-two classes is computed by
the operators $D_x+F(\lambda,L,\pi)$ in
\eqref{coh:lowering-k3} and \eqref{coh:lowering-torus}.

\begin{lemma}
 \label{ks:formal-kernel}
For the class $a$ defined by \eqref{ks:test-tensor}, its
annihilator $\mathcal Z\subset H^2(Y\times T_g,\Q)$ has
dimension $19$ and projects injectively to $H^2(Y,\Q)$.
In particular, it satisfies \eqref{setup:wedge-injective}.
\end{lemma}

\begin{proof}
Decompose
$x=x_N+x_e+x_0t_0$, where $x_e=\sum_sx_{e,s}e_s$.
The $e_s$ components of the annihilation equation are
\begin{equation}
 \label{ks:e-equations}
 F(R_s)+x_{e,s}U_-^*=0.
\end{equation}
For $s=2$, the second term has exterior degree two.  On a
power $\theta_i^k$, the contraction, derivation, and wedge
terms of $F$ have degrees $2k-2,2k,2k+2$, respectively.
For $k=4,7,10$ these nine degrees are distinct and exceed
two.  Hence $x_{e,2}=0$, and each of the nine terms vanishes
separately.

Symplectic Lefschetz injectivity now gives $\lambda=0$ from
$\lambda\wedge\theta^4=0$.  The derivation terms give
$F_L(\theta_i)=0$ for $i=0,1,2$, where $\theta_0=\theta$.
For the contraction term, contracting two $a$ directions into
$\theta^4$ gives, up to a nonzero factor, the two-$b$ tensor
obtained from $\pi$ by $Q$, wedged with $\theta^2$.
The same injectivity gives $\pi=0$.  All these wedge maps
are in the injective range for $g=2^{19}$.  The remaining
equations \eqref{ks:e-equations} give $x_e=0$.

The equation $F_L(\theta)=0$ says that $L$ is self-adjoint
for $Q$.  Preservation of $\theta_1$ and $\theta_2$ says
in addition that $L$ commutes with their relative operators.
If $x=0$, the $N$ equations become
$F_L(J(y))=0$ for every $y\in N$.  By the matrix description
above, these are exactly
$[L,\Gamma_y]=0$.  Irreducibility on $M_1$, followed by
the common-commutant condition on $R_0$, makes $L$ scalar.
The equation $F_L(U_0^*)=0$ kills this scalar, because
the lowering action on the highest-weight-two representation
is nonzero.  The kernel therefore projects injectively to
$N\oplus\Q t_0$, a space of dimension $19$.

Conversely, each $x\in N\oplus\Q t_0$ occurs in the kernel.
The operator $D_x$ is the lowering-degree element of
$\mathfrak{so}(V)$ displayed in \eqref{coh:lowering-k3}.
Its simultaneous spin action preserves \eqref{ks:test-tensor}.
Since $W$ has only Cartan weights $1$ and $-1$, its lowering
action is of the form $F_L$.  This proves the dimension
assertion.  Finally, projection of $H^3(Y\times T_g,\Q)$
to $H^2(Y,\Q)\otimes H^1(T_g,\Q)$ identifies the map in
\eqref{setup:wedge-injective} with the tensor product of
the injective projection of $\mathcal Z$ and the identity.
\end{proof}

\subsection{Determining the transcendental component}
\label{ks:exactness-section}

Proposition~\ref{prop:realization} now supplies a perfect
complex on $B\times E^g$.  Divide its Mukai vector by the
common nonzero rational scalar in the Euler comparison and
call the result $C$.  The dictionary determines
$C|_{P_B}=C_*|_{P_B}$, but does not initially determine
its component on $U_B=T(B)$.  The Ext bound does.  In the
following statement, give the formal space $U$ the Hodge
structure induced by the actual elliptic structure on $H$
through its highest-weight-two representation.

\begin{lemma}
 \label{ks:exactness}
There is a rational Hodge isometry $U_B\simeq U$, for the
actual Hodge structure on $H=H^1(E,\Q)$, under which
$C=C_*$.  The semiregularity trace of the perfect complex
is injective.
\end{lemma}

\begin{proof}
\emph{Excluding a zero transcendental component.}
Choose normalized actual Hodge directions $t_+,t_0,t_-$
on $U_B$, as in the cap calculation of
Section~\ref{sec:cohomology}, and write its unknown component
in the form \eqref{ks:weight-components}.
By Lemma~\ref{ks:formal-kernel} and the Ext bound, the actual
cap kernel has dimension at least $19$.  Imposing the three
linear conditions $x_e=0$ leaves a subspace of dimension at
least $16$.  On this subspace the
equations $F(R_s)=0$ again imply $\lambda=\pi=0$, and
make $L$ self-adjoint and commuting with the two relative
operators.  The $N$ equations are
\begin{equation}
 \label{ks:N-equation}
 F_L(J(y))+q(x_N,y)U_-=0\qquad(y\in N).
\end{equation}
The $t_-$ component, using \eqref{coh:lowering-k3}, is
\begin{equation}
 \label{ks:t-minus-equation}
 F_L(U_+)=J(x_N)+q(x_0t_0,t_0)U_0.
\end{equation}
If the entire component \eqref{ks:weight-components} were
zero, Equation~\eqref{ks:N-equation} would make $L$ scalar,
and \eqref{ks:t-minus-equation} would give $J(x_N)=0$,
hence $x_N=0$.
Only the scalar $L$ and $x_0$ would remain, contrary to
the dimension lower bound $16$.  Thus the component is
nonzero.

\emph{Identifying the Hodge representation.}
This nonzero component realizes the simple, three-dimensional
K3 Hodge structure $U_B$ in the exterior algebra of copies of
$H$, with the indicated Tate twists.  Since $E$ is non-CM,
its irreducible tensor representations are the irreducible
$\mathrm{SL}_2$ representations.  The only possible one here
has highest weight two.  The resulting rational Hodge
identification $U_B\simeq U$ is a similarity: the invariant
forms on this irreducible representation are proportional.
The underlying quadratic spaces are already rationally
isometric by the dictionary.  In odd dimension their similarity
multiplier is a rational square, as follows by taking
determinants; the signatures give the positive sign.  Rescaling
gives a Hodge isometry.  We may now use it to compare the two
sets of weight directions.  The component
\eqref{ks:weight-components} is equivariant for the simultaneous
$\mathrm{SL}_2$; in particular $U_-\ne0$.

\emph{Determining the component up to one scalar.}
Still on the subspace $x_e=0$, the kernel of projection
to $x_N$ has dimension at most two: setting $x_N=0$
in \eqref{ks:N-equation} again makes $L$ scalar, leaving
at most that scalar and $x_0$.  The image in $N_{\C}$
therefore has dimension at least $14$.  A totally isotropic
subspace of the $18$-dimensional nondegenerate space
$N_{\C}$ has dimension at most nine.  Choose a kernel
element for which $q(x_N,x_N)\ne0$.

For every $y\perp x_N$, Equation~\eqref{ks:N-equation}
says that $L$ commutes with $\Gamma_y$.  On the irreducible
Clifford module $M_1$, the commutant of these $17$
directions is spanned by $1$ and $\Delta\Gamma_{x_N}$;
this is the restriction of the even-dimensional Clifford
module to an odd-dimensional perpendicular.  The relative
operators again remove the multiplicity commutant, so
\[
 L=l\,1+m\,\Delta\Gamma_{x_N}.
\]
The Clifford relation gives
\[
 [\Delta\Gamma_{x_N},\Gamma_y]
   =\kappa\,q(x_N,y)\Delta
\]
for a fixed nonzero normalization constant $\kappa$.
Equation~\eqref{ks:N-equation} therefore makes $U_-$ a
nonzero multiple $cU_-^*$: the two-$b$ tensor with matrix
$\Delta Q$ is exactly the lowest-weight line of $J(U)$.
Equivariance gives this same scalar on all three
components in \eqref{ks:weight-components}.

\emph{Determining the exact normalization.}
Let $L^*$ be the spin lowering lift for
$x=x_N$ in the tensor $C_*$.  Comparing its version of
\eqref{ks:N-equation} with the actual equation gives
\[
 L=cL^*+\widetilde l\,1.
\]
Now use Equation~\eqref{ks:t-minus-equation} modulo the line
of $U_0$.  The scalar part $\widetilde l\,1$ of $L$
lowers $U_+$ into that line.  The equation for $C_*$,
with $x=x_N$, gives $F_{L^*}(U_+^*)=J(x_N)$.
Since $U_+=cU_+^*$, we obtain
\[
 J(x_N)\equiv F_L(U_+)
   \equiv c^2 F_{L^*}(U_+^*)=c^2J(x_N)
       \pmod{\C U_0}.
\]
One factor $c$ comes from $L$ and the other from $U_+$.
The tensor
$J(x_N)$ lies outside that line, since $J(N)$ and $J(U)$
have the distinct Clifford descriptions above.  Hence
$c^2=1$.  Reversing the sign of the transcendental Hodge
isometry if necessary gives $C=C_*$.

The circle on the resulting $V=h^\perp$ has its unique
lift to the spin group, and its action on $W$ is the
actual circle on $H^1(E^g,\Q)$.  Thus this actual
weight-one structure is the stipulated Kuga--Satake
structure.  The formal and actual cap calculations now
agree, so the cap kernel has dimension exactly $19$.
Lemma~\ref{coh:rank-trace} makes the
semiregularity trace injective.
\end{proof}

\subsection{Deformation and all polarization degrees}
\label{ks:deformation}

We have constructed an exact tensor whose complex can deform
whenever its Mukai vector stays Hodge.  We first use this to
move away from $B$ while preserving its polarization degree.
We then compare that degree with any prescribed degree by a
rational isometry, and specialize the exact tensor to fibers
with larger Picard rank.

\paragraph{Deformation from the mirror quartic.}
Consider the full family of $h$-polarized K3 surfaces, coupled
with their Kuga--Satake abelian varieties.  These are algebraic families
after a suitable level cover.  More precisely, the primitive
period domain and its spin cover, together with the invariant
rational alternating form, give a holomorphic equivariant
map to Siegel space.  Arithmetic subgroups under the
isogeny from Spin to $\SO$ have arithmetic image.  Intersect
with the stabilizers of the chosen lattices and pass to a
neat level; the K3 period theorem and algebraicity of the
arithmetic quotients and their morphisms then give a smooth
algebraic base with both universal families
\cite{DeligneK3,BailyBorel,BorelAlgebraicity}.

The lattice at the central point may be chosen to give
exactly the abelian variety $E^g$: preserve that lattice
by further level restriction and multiply the rational
polarization by an integer, with the sign chosen above.
On this spin level cover the tensors in
\eqref{ks:test-tensor} are global parallel sections:
the vectors $h$, $H^0$, and $H^4$ remain fixed and algebraic, the
sum on $V$ is Spin-invariant, and every $\theta_i$ is
Spin-invariant.  Hence the full Mukai vector $C_*$ stays
Hodge throughout the family; so do its Todd factors.
Lemma~\ref{ks:exactness} and
Proposition~\ref{prop:semiregularity} give algebraicity
of this class on every fiber of the polarized component.

On a fiber $S\times A$, take the component in $H^2(S)\otimes H^2(A)$
and project the K3 factor to
$V=h^\perp$.  The additional terms $e_s\otimes R_s$ have
abelian degrees $0$, $8$, $14$, $20$, or $26$, so these
algebraic projections extract $J$.  To regard $J$ as a correspondence
with two $H^1$ slots, use the surface and abelian
K\"unneth projectors, the polarization and its inverse,
and the inclusion of an exterior square into two ordered
slots.  All are algebraic: the inverse polarization is
provided by the Lefschetz standard conjecture for abelian
varieties, and the slot conversion can also be realized
using the addition map and pullback.

This already determines the usual standard tensor, not
only a tensor with the volume element on the right.
Since $z$ is central, $L_{vz}=L_vR_z$ on the even space.
For any nonisotropic $w\in V$, replace $R_z$ by $R_w$
using right multiplication by the invertible even
element $z^{-1}w$.  Every right even multiplication is
a weight-one Hodge endomorphism of the abelian variety,
and is therefore algebraic.  Thus the map
$x\mapsto vxw$ is algebraic in this full primitive
Clifford construction.

\paragraph{Comparison of polarization degrees.}
Fix a prescribed degree $2d$.  The initial rational ample
class $h$ on $B$ can be chosen with square $2d$.  Indeed,
Meyer's theorem gives an isotropic rational vector in the
indefinite rank-$19$ space $\NS(B)_{\Q}$, hence a rational
hyperbolic plane.  The quadric $q(x,x)=2d$ consequently has
rational points, and rational approximation on this quadric
gives a point in the open ample cone.  Scaling to a primitive
integral class changes its square by a rational square.
Thus the primitive integral class chosen on $B$ has square $2d$
times a rational square.  Witt's extension theorem
therefore gives a rational isometry of the K3 quadratic
spaces carrying its polarization line to the line of
a primitive vector of square $2d$.  The isometry
identifies the primitive period domains.  At very general
periods on both sides the surfaces have Picard rank one:
the excluded wall hyperplanes, including their translates
under the rational isometry, are rational hyperplanes
and can all be avoided.  The K3 period theorem supplies
both polarized surfaces.  Buskin's theorem makes their
rational Hodge isometry algebraic \cite{Buskin}.
If $f$ denotes its restriction to the primitive spaces,
the induced Clifford Hodge isomorphism satisfies
\[
 C^+(f)\,j_w(v)\,C^+(f)^{-1}
       =j_{f(w)}(f(v)).
\]
This weight-one isomorphism is induced by an abelian
isogeny up to a rational multiple.  Composing it with
the K3 correspondence, and converting polarizations
as above, therefore transports the exact tensor.
We obtain the exact standard tensor very generally in
degree $2d$.  One may make this comparison on marked
components or on commensurable level quotients of the
two domains.

\paragraph{Specialization and transcendental restriction.}
Specialization gives the full primitive tensor at
every polarized fiber.  Here the class being specialized
is the prescribed parallel tensor in the full
Kuga--Satake family, so Picard jumps cause no loss of
information.  Equivalently, apply the specialization
theorem and its primitive-map corollary in
\cite[Theorem~1.2 and Corollary~5.5]{CKS}: the very-general nonzero
algebraic input required there has just been supplied
in its stronger exact form.  This specialization is
the proper-cycle-parameter argument of
Proposition~\ref{prop:semiregularity}; it does not
replace actual algebraicity by the Hodge condition.

\begin{proof}[Proof of Theorem~\ref{thm:exact-ks}]
Every projective K3 surface has a primitive ample class,
so the preceding construction gives its full primitive
Kuga--Satake tensor.  The projection to $T(S)$ is algebraic,
using the divisor intersection pairing and the degree-zero
and top-degree projectors.  For the prescribed
$w\in T(S)$, restrict the even Clifford space to
\[
 C^+(T(S))\subset C^+(h^\perp).
\]
This is a rational weight-one sub-Hodge structure, and
$x\mapsto vxw$ preserves it for $v\in T(S)$.  It can
therefore be split off by algebraic abelian
correspondences.  Changing its lattice, isogenous
realization, or rational polarization also uses algebraic
abelian correspondences.  The resulting cycle has the
exact action \eqref{ks:standard-map} with the chosen
polarization conversion, proving the theorem.
\end{proof}

\begin{corollary}
 \label{cor:self-powers}
For every projective complex K3 surface $S$, the rational
Hodge conjecture holds in every degree on every self-power
$S^m$.
\end{corollary}

\begin{proof}
Theorem~\ref{thm:exact-ks} supplies the exact standard
tensor on the full even Clifford space, for one
nonisotropic $w$ and one polarization.  These are the
hypotheses of the general self-power criterion
\cite[Theorem~1.1]{AKS}.  That criterion is independent
of the special quadratic locus used elsewhere in the
same reference to construct correspondences.
\end{proof}
